%% ma: L12-B08-0097
%% lop: 12
%% chuong: 2
%% bai: 8
%% dang: 12.8.29
%% loai: TN4
%% muc: TH
%% dap_an: "C"
%% nguon: "CĐ2.3.pdf giáo viên gửi 27/10/2026 (M07, Lê Minh Khôi – Nguyễn Anh Khôi, Chương II Chủ đề 3 Biểu thức tọa độ của các phép toán vectơ), Câu 6"
%% kien_thuc: "Bài 8 (biểu thức tọa độ của phép cộng, trừ vectơ, tích vectơ với một số), Bài 7 (tọa độ vectơ theo i, j, k)"
%% trang_thai: chinh-thuc
\begin{ex}
Trong không gian $Oxyz$, cho hai vectơ $\vec a=2\vec i+3\vec j$, $\vec b=(9;1;-2)$ và hai điểm $A(2;0;3)$, $B(6;7;-1)$. Tọa độ của vectơ $\vec u=3\cdot\vec a-\vec b+\overrightarrow{AB}$ là
\choice{$\vec u=(-1;-15;2)$}{$\vec u=(7;-1;-6)$}{\True $\vec u=(1;15;-2)$}{$\vec u=(-7;1;6)$}
\loigiai{$\vec a=(2;3;0)$, $\overrightarrow{AB}=(4;7;-4)$.

$\vec u=(6-9+4;9-1+7;0+2-4)=(1;15;-2)$.}
\end{ex}
